\section{Three expiries and a broken ruler}
\label{sec:clkboard}

Part~II has been rebuilding the observation, one input at a time.
Chapter~6 fitted the forward and the discount; Chapter~7 removed the
early-exercise right and handed every quote back as a European price.
One input has ridden along unexamined since Chapter~2: the maturity
variable $\tau$ that sits under every total variance $w=\sigma^2\tau$.
The book has been calling it \emph{variance time} from the start ---
while quietly setting it equal to the calendar, since no chapter yet
had a reason to distinguish them.  Chapter~7 closed by naming the
debt: ``two clocks this book has so far been allowed to conflate,
because its frozen data happens to agree on both.''  This chapter is
about the days when they disagree.

Start with real prices, because the disagreement is sitting in the
book's own frozen data.  The snapshot behind every chapter --- the
Monday evening of 2026-08-03 --- carries two names, and so far the
chapters have worked mostly on SPY.  Turn instead to NVDA, the
snapshot's single-stock name, and read its at-the-money implied
volatility across the eight quoted expiries --- the readings computed
as $\sqrt{w/t}$, total variance over calendar time, which this
chapter will call the \emph{calendar reading} and write
$\sigma_{\rm cal}$.  (At the money means at $k=0$, the forward; the
precise read-off rule is one line in
appendix~\ref{app:clkprotocol}, which also lists the eight expiry
dates.  As before, one \emph{vol bp} is $10^{-4}$ of absolute
volatility; one \emph{variance bp} is $10^{-4}$ of total variance.)
Three of the eight readings make the puzzle by themselves:
\[
\underbrace{\MacClkBoardVolFourDPct\%}_{\text{4 days}}
\qquad
\underbrace{\MacClkBoardVolEighteenDPct\%}_{\text{18 days}}
\qquad
\underbrace{\MacClkBoardVolFortySixDPct\%}_{\text{46 days}} .
\]
The same underlier, the same afternoon, three expiries two and four
weeks apart --- and the middle one is quoted
\MacClkBoardDipBp\ vol bp below the shortest.  Is NVDA's volatility
falling or rising with maturity?  The honest answer is that the
question is broken, and the goal of this section is to see why.

Ask the question in the right units.  A volatility is a rate; before
trusting a rate, look at the quantity it is the rate \emph{of}.  The
quantity is total implied variance $w=\sigma^2 t$ --- for now we write
$t$, the calendar year fraction of Chapter~6, because the calendar is
the only clock we have so far.  The three expiries carry
\[
w(\text{4d})=\MacClkBoardWFourBp,\qquad
w(\text{18d})=\MacClkBoardWEighteenBp,\qquad
w(\text{46d})=\MacClkBoardWFortySixBp
\quad\text{variance bp.}
\]
Now do the division a first course would do: variance gained per
calendar day, interval by interval.  Between the 4-day and the 18-day
expiries, \MacClkBoardLullDays\ calendar days pass, so
\[
\frac{\MacClkBoardWEighteenBp-\MacClkBoardWFourBp}
     {\MacClkBoardLullDays}
=\MacClkBoardLullPerDay\ \text{variance bp per day;}
\]
between the 18-day and the 46-day expiries, \MacClkBoardHotDays\ days
pass, and
\[
\frac{\MacClkBoardWFortySixBp-\MacClkBoardWEighteenBp}
     {\MacClkBoardHotDays}
=\MacClkBoardHotPerDay\ \text{variance bp per day.}
\]
The days are not alike.  A day in the second window carries
\MacClkBoardHotOverLull\ times the variance of a day in the first.
And the second window is not anonymous.  The 18-day expiry is
August~21; the 46-day expiry is September~18; and NVDA reports
earnings once a quarter, on a schedule published well in advance,
with the next report due in late August --- after the first of those
dates, before the second.  The market is quoting exactly what it
believes: the report's day will move the stock like several ordinary
days, every expiry that \emph{contains} that day inherits its
variance, and the expiry that \emph{dodges} it --- August~21 --- gets
to be cheap.

\Cref{fig:clkboard} shows the whole board.  In panel~(a), the orange
curve is the calendar reading $\sigma_{\rm cal}$ of all eight
expiries against calendar days to expiry, on a logarithmic axis.
Notice the shape around the annotated dip: the two shortest expiries
(2 and 4 days, both before August~21) read above $43\%$, the 18-day
expiry reads \MacClkBoardVolEighteenDPct\%, and the 46-day expiry
recovers most of the drop, at \MacClkBoardVolFortySixDPct\%.
Further out, the curve wanders rather than trends --- the highest
reading on the whole board belongs to the 228-day expiry.  Read as
``the market's forecast of NVDA's volatility as a function of
horizon,'' the curve is erratic.  Panel~(b) draws the same eight
quotes as what they actually pin down: accrual rates
$\Delta w_i/\Delta t_i$, as horizontal steps --- one step per
interval, the first interval running from the valuation date (where
accrued variance is zero) to the first expiry, each later one
between consecutive expiries.  (This per-interval rate is called
\emph{forward variance}, and \cref{sec:clkinverse} gives it a symbol
and a starring role.)  The two computations above are the third and
fourth steps: the low step is the pre-report lull, and the shaded
interval --- the one containing the scheduled report --- runs hot.
Panel~(b) is panel~(a) with the division done honestly: nothing here
is erratic, it is one hot window sitting among cooler ones.

\begin{figure}[htbp]
\centering
\paperfig{\textwidth}{fig_clk_board.pdf}
\caption{The opening puzzle (NVDA, frozen snapshot of 2026-08-03).
(a)~At-the-money implied volatility against calendar days to expiry:
the 18-day expiry --- the one that expires \emph{before} the scheduled
late-August earnings report --- is quoted \MacClkBoardDipBp\ vol bp
below the 4-day reading.  (b)~The same board as accrual rates: total
variance gained per calendar year in each interval between consecutive
expiries.  The interval that contains the report (shaded) accrues
\MacClkBoardHotOverLull$\times$ the variance per day of the lull
interval before it.}
\label{fig:clkboard}
\end{figure}

Why call this a \emph{broken ruler} rather than an interesting market
fact?  Because of what it does to every tool the book has built.  Each
single-expiry fit of Part~I is untouched --- a smile model fits its own
expiry's quotes and never asks how fast variance accrued on the way ---
but everything that couples expiries inherits the kink.  A
term-structure view drawn through these readings shows volatility
falling four points in two weeks and then climbing back, which no
trader believes is a forecast.  An interpolation between the quoted
expiries must invent readings for unlisted maturities, and
\cref{sec:clkinterp} shows the calendar-clock rule inventing them
badly.  The local volatility field of Chapter~4 consumes
$\partial_\tau w$ directly --- its numerator \emph{is} the accrual
rate of panel~(b).  Even the calendar-order condition of Chapter~2
(total variance non-decreasing in maturity at fixed $k$) offers no
comfort: it holds on this board, and the board is still a mess to
read, because order says nothing about the accrual being bunched
into one window.

There are two ways to respond.  One is to enrich the model: add a jump
of random size at the report date, price options on the mixture, and
carry the machinery of jump models from here on.  That road is real
but expensive, and this book does not take it (\cref{sec:clkfrozen}
says where it lives in the literature).  The other response is older
and much cheaper: keep the model, change the \emph{clock}.  If an
earnings day is worth several ordinary days of variance, then measure
time in units in which it \emph{is} several days.  On that clock,
variance accrues at a steady rate again, and the constant-rate
diffusion picture --- the one behind every formula in this book ---
fits without a fight.  The next section shows that this is not a
trick: it is how the mathematics of martingales says time should be
measured in the first place.
